% Sommateur : cercle barré d'une croix, signes à l'intérieur de chaque quadrant
% (#1 haut, #2 droite, #3 bas, #4 gauche). Taille fixe en pt pour que les signes
% restent dans le cercle quelle que soit la police du schéma.
\tikzset{add/.style n args={4}{
    circle, minimum size=13pt, inner sep=0pt,
    append after command={
        \pgfextra{\let\jvsomme\tikzlastnode}
        ($(\jvsomme.center)+(45:6.5pt)$)  edge[draw=black, -] ($(\jvsomme.center)+(225:6.5pt)$)
        ($(\jvsomme.center)+(135:6.5pt)$) edge[draw=black, -] ($(\jvsomme.center)+(315:6.5pt)$)
        ($(\jvsomme.center)+(0,4.1pt)$)  node[jvsigne] {#1}
        ($(\jvsomme.center)+(4.1pt,0)$)  node[jvsigne] {#2}
        ($(\jvsomme.center)+(0,-4.1pt)$) node[jvsigne] {#3}
        ($(\jvsomme.center)+(-4.1pt,0)$) node[jvsigne] {#4}
    }
}}
\tikzset{jvsigne/.style={rectangle, draw=none, fill=none, minimum size=0pt,
    inner sep=0pt, text=black, font=\fontsize{4pt}{4pt}\selectfont}}

\tikzstyle{block} = [draw, fill=white, rectangle, align=center, font=\small]%
\tikzstyle{comp} = [draw, fill=white, circle, add={}{}{$-$}{$+$}, node distance=0.5cm]%
\tikzstyle{sum3} = [draw, fill=white, circle, add={$+$}{}{$+$}{$+$}, node distance=0.5cm]%
\tikzstyle{sum2} = [draw, fill=white, circle, add={$+$}{}{}{$+$}, node distance=0.5cm]%
\tikzstyle{input} = [text centered, inner sep=0pt, outer sep=0pt, font=\footnotesize]%
\tikzstyle{output} = [text centered, inner sep=0pt, outer sep=0pt, font=\footnotesize]%

\begin{tikzpicture}[auto,node distance=1cm,>=latex']
    \footnotesize
	% --------------------------------------
	% --- définition des espacements ---
	% --------------------------------------
	\def\x{2.1em}
	\def\y{1.8em}
	\def\r{3pt} % rayon horizontal du pont
	\def\h{1.3pt} % hauteur du pont

	% ------------------------------
	% --- définition des nodes ---
	% ------------------------------
	% node de consignes
	\node [input] (r2) {$r_2(k)$};

    % node du préfiltre
    \node [block, anchor=north west, blocRegul] (PF) at ([xshift=\x, yshift=-\y]r2.east) {$P\!F(q)$};

    % nodes des comparateurs
    \node [comp, anchor=west, blocRegul] (compP) at ([xshift=\x]PF.east|-r2) {};
    \node [comp, anchor=west, blocRegul] (compI) at ([xshift=\x]PF.east-|compP) {};

    % nodes régulateurs
    \node [block, anchor=west, blocRegul] (Ci) at ([xshift=1.6*\x]compI.east) {$C_I(q)$};
    \node [block, blocRegul] (Cp) at (compP.east-|Ci) {$C_P(q)$};

    % node sommateur
    \node [sum3, anchor=west, blocRegul] (sumCU) at ([xshift=\x]Cp.east) {};

    % node de la commande externe
    \node [input, anchor=south] (r1) at ([yshift=\y]sumCU.north) {$r_1(k)$};

    % node procédé
    \node [block, anchor=west, blocProcede] (Go) at ([xshift=\x]sumCU.east) {$G_o(q)$};

    % node sommateur du bruit
    \node [sum2, anchor=west, blocBruit] (sumH) at ([xshift=\x]Go.east) {};

    % node modèle de bruit
    \node [block, anchor=south, blocBruit] (Ho) at ([yshift=\y]sumH.north) {$H_o(q)$};

    % node du bruit 
    \node [input, anchor=east] (eo) at ([xshift=-\x]Ho.west) {$e_o(k)$};

    % node de la sortie
    \node [input, anchor=west] (y) at ([xshift=\x]sumH.east) {$y(k)$};

	% -------------------------------
	% --- définition des flèches ---
	% -------------------------------
	% de r2 aux comparateurs en passant par PF
	\draw [->] (r2) -- (compP);
	\draw[->] ($(r2.east)+(0.5*\x,0)$) |- (PF);
	\path[name path=pathPFCI] (PF) -- (compI);

	% des comparateurs aux régulateurs
	\draw [->] (compP) -- (Cp) node [midway] {$\varepsilon_P(k)$};
	\draw [->] (compI) -- (Ci) node [midway] {$\varepsilon_I(k)$};

	% des régulateurs au sommateur
	\draw [->] (Cp) -- (sumCU);
	\draw [->] (Ci) -| (sumCU);
	\draw [->] (r1) -- (sumCU);

    % du sommateur au procédé
    \draw [->] (sumCU) -- (Go) node [midway] {$u(k)$};

    % du procédé au sommateur 
    \draw [->] (Go) -- (sumH);

    % du modèle de bruit au sommateur
    \draw [->] (Ho) -- (sumH) node[midway] {$\xi_0(k)$};

    % de l'entré du bruit au modèle de bruit
    \draw [->] (eo) -- (Ho);

    % du sommateur à la sortie
    \draw [->] (sumH) -- (y);

    % du sommateur aux compateurs 
    \draw [->] (sumH) -| ([xshift=\x/2, yshift=-\y]Ci.south-|sumH.east) -| (compI);
    \draw [->, name path=pathSFCP] (sumH) -| ([xshift=\x/2, yshift=-\y]Ci.south-|sumH.east) -| (compP);

    % calcul de l'intersection entre le préfiltre--comparateur intégral et différenciateur--comparateur proportionnel
	\path[name intersections={of=pathPFCI and pathSFCP, by=crossPFI}];
	\draw[->] (PF) -- ($(crossPFI)+(-\r,0)$) arc[start angle=180,end angle=0,x radius=\r,y radius=\h] -- (compI);
\end{tikzpicture}